\documentclass[hidelinks]{article}
\usepackage{Resources/UoBLab}
\usepackage[spanish]{babel}
\usepackage[utf8]{inputenc}
\usepackage{amsmath}
\usepackage{amssymb}
\usepackage{physics}

\setlength{\parskip}{0pt}

\pubyear{2026}
\subjectarea{}

\begin{document}

\title{Examen 1}
\author{Edgar Agustin Contreras}
\course{Física del Estado Sólido}
\school{Centro de Investigación y de Estudios Avanzados del IPN (CINVESTAV)}
\date{2 de octubre de 2026}
\keywords{estructura de bandas, modelo de Kronig--Penney, masa efectiva}

\maketitle

\setcounter{tocdepth}{2}


\section{Introducción}

Hoy en día, el desarrollo de dispositivos electrónicos, como los emisores de luz (LEDs), los transistores, los láseres y los paneles solares, está basado en los semiconductores. Estos materiales se distinguen por una brecha de energía entre la banda de valencia y la de conducción. Ante una perturbación, como un campo eléctrico o un esfuerzo mecánico, sus propiedades ópticas y eléctricas se modifican.

La herramienta más usada para describir ese comportamiento es la estructura de bandas, que relaciona la energía de los electrones con su vector de onda cristalino $\mathbf{k}$. El análisis y el cálculo de esta relación permiten entender el material, por lo que el modelado resulta importante.

El modelado de un sistema de muchas partículas es, sin embargo, un problema muy complejo, así que hace falta introducir simplificaciones. Una consiste en colocar los átomos en un arreglo periódico, descrito por una red de Bravais. Otra simplifica el problema a los electrones de valencia en el potencial de los iones: esos electrones son, en general, los responsables de las propiedades electrónicas del sólido.

En este trabajo se calcula la estructura de bandas de una cadena monoatómica de aluminio metálico orientada en la dirección $[111]$, esto es, la dirección que atraviesa los planos de mayor empaquetamiento, los planos $(111)$. El cálculo usa el modelo de Kronig--Penney y, como casos particulares, el electrón libre y el electrón casi libre.

\section{Metodología}

\subsection{Códigos}

Las simulaciones se hicieron en MATLAB R2024b. El cálculo se reparte en cinco archivos, que se reproducen en el apéndice:
\begin{itemize}
\item \texttt{main.mlx}: \emph{live script} que fija los parámetros y, en una sección por cada caso estudiado, grafica la estructura de bandas llamando a los demás archivos.
\item \texttt{kv.m}: evalúa la relación de dispersión~\eqref{eq:dispersion} y devuelve $k=\arccos(F)/a$.
\item \texttt{k.m}: recibe un arreglo de energías y devuelve la estructura de bandas, separando la parte real de $k$ de la imaginaria.
\item \texttt{fek.m}: devuelve $k(E)$ del electrón libre, ecuación~\eqref{eq:libre}.
\item \texttt{bordes.m}: localiza los bordes de banda en el arreglo que devuelve \texttt{k.m} y de ahí la anchura de la brecha $E_g$.
\end{itemize}

Las longitudes se expresan en ångströms, el vector de onda en \AA$^{-1}$ y la energía en eV. Con ese sistema de unidades, la masa del electrón y $\hbar$ se introducen con los valores numéricos de la práctica, de modo que $k=\sqrt{2m_e E}/\hbar$ queda directamente en \AA$^{-1}$. El periodo del modelo no es la constante de red cúbica del aluminio, $a_{\mathrm{Al}}=4.05$~\AA, sino la distancia entre planos $(111)$:
\begin{equation}
a=\frac{a_{\mathrm{Al}}}{\sqrt{3}}=2.3383~\text{\AA}.
\label{eq:periodo}
\end{equation}

\subsection{Modelo de Kronig--Penney}

El modelo de Kronig--Penney trata un solo electrón en un cristal unidimensional \cite{kronig1931}. El potencial de los iones se sustituye por un potencial periódico y constante a trozos, lo que permite una solución analítica. En cada periodo hay dos regiones, como se muestra en la figura~\ref{fig:kp}. La región~II coincide con los iones y la región~I con el espacio entre ellos. El periodo es $a=d+s$.

\begin{figure}[htbp]
\centering
\IfFileExists{../../../../Pictures/Screenshots/Screenshot 2026-10-01 220916.png}{%
  \includegraphics[width=0.72\linewidth]{../../../../Pictures/Screenshots/Screenshot 2026-10-01 220916.png}%
}{%
  \fbox{\parbox{0.72\linewidth}{\centering\vspace{2.5em}
  Potencial periódico de Kronig--Penney\vspace{2.5em}}}%
}
\caption{Potencial del modelo de Kronig--Penney. La región~I, entre iones, es una barrera de altura $V_0$ y ancho $s$. La región~II, sobre los iones, es un pozo de potencial nulo y ancho $d$. El periodo es $a=d+s$.}
\label{fig:kp}
\end{figure}

La función de onda satisface la ecuación de Schrödinger
\begin{equation}
H\psi=
\left(-\frac{\hbar^2}{2m_e}\frac{d^2}{dx^2}+V_{\mathrm{KP}}(x)\right)\psi
=E\psi.
\label{eq:schrodinger}
\end{equation}
El cero de la energía es arbitrario. Con el origen de la figura~\ref{fig:kp}, el potencial de un periodo se escribe
\begin{equation}
V_{\mathrm{KP}}(x)=
\begin{cases}
V_0, & an \leq x \leq an+s, \quad n\in\mathbb{Z},\\
0, & \text{en el resto del periodo,}
\end{cases}
\label{eq:potencial}
\end{equation}
donde el ancho de la barrera es $s=a-d$ y el del pozo es $d$. La región~I es la barrera y la región~II es el pozo.

En cada región la solución general es una combinación de ondas planas. Para $E>V_0$,
\begin{equation}
\begin{aligned}
\psi_{\mathrm{I}}(x)
&=A\exp[i\beta x]+B\exp[-i\beta x],
&\beta&=\frac{\sqrt{2m_e(E-V_0)}}{\hbar},\\[4pt]
\psi_{\mathrm{II}}(x)
&=C\exp[i\alpha x]+D\exp[-i\alpha x],
&\alpha&=\frac{\sqrt{2m_e E}}{\hbar}.
\end{aligned}
\label{eq:ondas}
\end{equation}
Si $E<V_0$, $\beta$ es imaginario: en la barrera las ondas pasan a ser exponenciales reales. La relación de dispersión que se obtiene más abajo sigue siendo válida en ese caso.


Como el potencial es periódico, el teorema de Bloch escribe cada estado estacionario en la forma $\psi_k(x)=u_k(x)\exp(ikx)$, con $u_k(x+a)=u_k(x)$. Imponer eso junto con la continuidad de $\psi$ y de $d\psi/dx$ en la interfaz entre las regiones~I y~II deja un sistema homogéneo para $A$, $B$, $C$ y $D$, que tiene solución distinta de la trivial solo si
\begin{equation}
\begin{aligned}
\cos(ka)&=F,\\
F&=\cos(\alpha d)\cos(\beta s)
-\frac{\alpha^2+\beta^2}{2\alpha\beta}\sin(\alpha d)\sin(\beta s).
\end{aligned}
\label{eq:dispersion}
\end{equation}
$E$ entra en~\eqref{eq:dispersion} a través de $\alpha$ y $\beta$.

Cuando $|F|\leq 1$ existe un $k$ real y el estado es una onda de Bloch que se propaga por la cadena. En la primera zona de Brillouin, $-\pi/a\leq k\leq\pi/a$. El borde de banda es $|F|=1$, es decir $k=0$ o $k=\pm\pi/a$. Cuando $|F|>1$ no hay $k$ real: no hay electrones de Bloch propagantes y ese intervalo de energía es una brecha del modelo. El vector de onda es entonces imaginario.

\subsection{Electrón libre y electrón casi libre}

En el modelo de electrón libre los electrones no sienten el potencial de los iones. Dentro de Kronig--Penney ese caso se obtiene con $V_0=0$ y, como primera elección, $d=a/2$. El hamiltoniano se reduce a la energía cinética y la energía de las ondas planas es
\begin{equation}
E(k)=\frac{\hbar^2 k^2}{2m_e}.
\label{eq:libre}
\end{equation}
La curva~\eqref{eq:libre} es la referencia con la que se compara la banda que devuelve la ecuación~\eqref{eq:dispersion} al apagar la barrera.

Si el potencial de los iones es una perturbación débil, el mismo potencial se trata en la aproximación de electrón casi libre. La barrera es baja cuando $V_0$ es mucho menor que la energía cinética en el borde de la primera zona, que con el periodo~\eqref{eq:periodo} vale $\hbar^2(\pi/a)^2/(2m_e)=6.88$~eV. En esta práctica $V_0=0.1$~eV, unas setenta veces menor.

Por la periodicidad de la red, el potencial admite una serie de Fourier de periodo $a$. En la base de ondas planas del electrón libre, $V_{\mathrm{KP}}$ solo acopla estados cuyos vectores de onda difieren en un múltiplo de $2\pi/a$. En la frontera de la zona, $k=\pi/a$, ese acoplamiento abre la primera brecha, de anchura aproximada $2|V_1|$ \cite{kittel}. El primer coeficiente, que se calcula en el apéndice, es
\begin{equation}
|V_1|=\frac{V_0}{\pi}\left|\sin\!\left(\frac{\pi d}{a}\right)\right|,
\label{eq:V1}
\end{equation}
proporcional a $V_0$. Con $d=a/2$ se reduce a $2|V_1|=2V_0/\pi$.



\section{Resultados}

\subsection{Electrón libre}

La ecuación~\eqref{eq:dispersion} no da $E(k)$ en forma explícita, sino lo contrario: se resuelve como $k=\arccos(F)/a$, con $F=F(E)$. La variable independiente del cálculo es entonces la energía y no el vector de onda, así que las bandas se obtienen evaluando $k$ sobre una malla de energías. Para el caso del electrón libre se recorrió $E$ de $0.1$ a $30$~eV en pasos de $0.01$~eV, es decir 2991 puntos.

El límite inferior no puede ser $E=0$, porque ahí $\alpha=0$ y el factor $(\alpha^2+\beta^2)/(2\alpha\beta)$ de la ecuación~\eqref{eq:dispersion} queda indeterminado. Se tomó $0.1$~eV, que corresponde a $k=0.162$~\AA$^{-1}$, apenas un $12\%$ del borde de zona $\pi/a=1.3436$~\AA$^{-1}$, de modo que no se pierde información del fondo de la banda. El límite superior se eligió para cubrir las dos primeras zonas: el borde de la primera está en $E(\pi/a)=6.88$~eV y el de la segunda en $E(2\pi/a)=27.51$~eV, así que a $30$~eV la parábola libre llega a $k=2.81$~\AA$^{-1}\approx 2.09\,\pi/a$ y alcanza a plegarse una vez completa dentro de la primera zona. El paso de $0.01$~eV solo controla la suavidad de la curva, porque con $V_0=0$ no hay brechas que resolver.

Por último, $\arccos$ devuelve valores en $[0,\pi]$, de modo que la ecuación~\eqref{eq:dispersion} solo entrega $k\geq 0$. Como $E(k)=E(-k)$, la rama negativa se construye por reflexión de la positiva, y entre las dos se inserta un \texttt{NaN} para que la gráfica no las una con un segmento horizontal que no es parte de la solución.

La comparación se muestra en la figura~\ref{fig:libre}. La línea azul es la estructura de bandas de Kronig--Penney con $V_0=0$, calculada con \texttt{k.m}, y la línea roja punteada es el electrón libre exacto, ecuación~\eqref{eq:libre}, calculado con \texttt{fek.m}. Las dos curvas están superpuestas hasta $k=\pi/a=1.34$~\AA$^{-1}$. A partir de ahí el electrón libre sigue creciendo monótonamente con $|k|$, mientras que Kronig--Penney devuelve siempre un $k$ dentro de $-\pi/a\leq k\leq\pi/a$: es el esquema de zona reducida, en el que la parábola libre aparece plegada. Al desplegarla al esquema extendido ambas curvas se superponen, es decir, los dos modelos dan la misma relación de dispersión. No hay ningún intervalo de energía sin un $k$ real asociado, o sea que no hay brecha; era de esperarse, porque al apagar la barrera el potencial no tiene componentes de Fourier que acoplen los estados $k$ y $k-2\pi/a$.

\begin{figure}[htbp]
\centering
\IfFileExists{figuras/libre.png}{%
  \includegraphics[width=0.72\linewidth]{figuras/libre.png}%
}{%
  \fbox{\parbox{0.72\linewidth}{\centering\vspace{2.5em}
  Estructura de bandas con $V_0=0$\vspace{2.5em}}}%
}
\caption{Estructura de bandas del Al en la dirección $[111]$ con $V_0=0$. Azul: modelo de Kronig--Penney, ecuación~\eqref{eq:dispersion}, en el esquema de zona reducida. Rojo punteado: electrón libre exacto, ecuación~\eqref{eq:libre}.}
\label{fig:libre}
\end{figure}



\section{Discusión}



\section{Conclusiones}





% El entorno thebibliography ya pone el encabezado "Referencias" por si solo:
% agregar \section{Referencias} aqui lo duplicaba.
\begin{thebibliography}{99}
\bibitem[Kronig y Penney(1931)]{kronig1931}
R.~de~L.~Kronig and W.~G.~Penney,
Proc.~R.~Soc.~Lond.~A \textbf{130}, 499 (1931).

\bibitem[Kittel(1996)]{kittel}
C.~Kittel, \emph{Introduction to Solid State Physics}, 7.\textsuperscript{a}~ed.
(John Wiley \& Sons, New York, 1996).
\end{thebibliography}




\section{Apéndice}

\subsection{Fourier}

\noindent Podemos definir el potencial de Kronig-Penney:
\begin{equation*}
V(x) = \begin{cases} 
V_0, & 0 \leq x \leq s = a-d \\ 
0, & a-d < x \leq a 
\end{cases}
\end{equation*}

\noindent La componente de Fourier se define como:
\begin{equation*}
V_1 = \frac{1}{a} \int_{c}^{c+a} V(x) \exp\left(-\frac{i 2\pi x}{a}\right) dx
\end{equation*}

\noindent Con $c=0$ y partiendo con la definición anterior:
\begin{equation*}
V_1 = \frac{1}{a} \int_{0}^{a-d} V_0 \exp\left(-\frac{i 2\pi}{a}x\right) dx
\end{equation*}

\noindent Haciendo el cambio de variable:
\begin{align*}
u &= -\frac{i 2\pi}{a}x \\
u(0) &= 0 \\
u(a-d) &= -\frac{i 2\pi}{a}(a-d) = -i 2\pi \left(1 - \frac{d}{a}\right)
\end{align*}

\noindent Resolviendo la integral:
\begin{align*}
V_1 &= i \frac{V_0}{2\pi} \int_{0}^{-i 2\pi\left(1 - \frac{d}{a}\right)} \exp(u) du \\
&= i \frac{V_0}{2\pi} \exp(u) \bigg|_{0}^{-i 2\pi\left(1 - \frac{d}{a}\right)} \\
&= i \frac{V_0}{2\pi} \left[ \exp(-i 2\pi) \exp\left(i 2\pi \frac{d}{a}\right) - 1 \right] \\
&= \frac{i V_0}{2\pi} \left[ \cos\left(2\pi \frac{d}{a}\right) - 1 + i \sin\left(2\pi \frac{d}{a}\right) \right] \\
&= \frac{V_0}{2\pi} \left[ -\sin\left(2\pi \frac{d}{a}\right) + i \left(\cos\left(2\pi \frac{d}{a}\right) - 1\right) \right]
\end{align*}

\noindent Luego, calculamos la magnitud $|V_1|$:
\begin{align*}
|V_1| &= \frac{V_0}{2\pi} \left[ \sin^2\left(2\pi \frac{d}{a}\right) + \left(\cos\left(2\pi \frac{d}{a}\right) - 1\right)^2 \right]^{1/2} \\
&= \frac{V_0}{2\pi} \left[ \sin^2\left(2\pi \frac{d}{a}\right) + \cos^2\left(2\pi \frac{d}{a}\right) + 1 - 2\cos\left(2\pi \frac{d}{a}\right) \right]^{1/2} \\
&= \frac{V_0}{2\pi} \left[ 2 - 2\cos\left(2\pi \frac{d}{a}\right) \right]^{1/2} \\
&= \frac{V_0}{\pi} \left[ \frac{1 - \cos\left(2\pi \frac{d}{a}\right)}{2} \right]^{1/2}
\end{align*}

\noindent Como:
\begin{equation*}
\cos(2x) = \cos^2 x - \sin^2 x = 1 - 2\sin^2 x
\end{equation*}

\noindent Finalmente:
\begin{equation*}
\sin x = \sqrt{\frac{1 - \cos(2x)}{2}}
\end{equation*}

\begin{equation*}
|V_1| = \frac{V_0}{\pi} \left| \sin\left(\pi \frac{d}{a}\right) \right|
\end{equation*}

\vspace{0.5cm}
\noindent Demostrando que en efecto es proporcional a $V_0$.


\subsection{Códigos}

% Para incluir los listados hace falta \usepackage{listings} en el preambulo
% y luego, por cada archivo:
%   \lstinputlisting[language=Matlab,caption={kv.m}]{kv.m}

\noindent Aquí van los listados de \texttt{main.mlx}, \texttt{k.m}, \texttt{kv.m}, \texttt{fek.m} y \texttt{bordes.m}.







\end{document}
